\documentclass[aspectratio=169,11pt]{beamer}

\usepackage[UTF8]{ctex}
\usepackage{amsmath,amssymb,booktabs,graphicx,listings,mwe}

\usetheme[footline=full]{USTS}

\hypersetup{hidelinks}
\graphicspath{{assets/}}

\lstset{
  basicstyle=\ttfamily\small,
  keywordstyle=\color{USTSDeepBlue}\bfseries,
  commentstyle=\color{USTSGray},
  frame=single,
  rulecolor=\color{USTSLightBlue},
  breaklines=true,
  columns=fullflexible,
  keepspaces=true,
  showstringspaces=false
}

\title[USTS Beamer]{苏州科技大学 Beamer 通用模板}
\subtitle{适用于课程汇报、学术报告与项目展示}
\author[绿化三]{绿化三}
\institute[USTS]{物理科学与技术学院\\苏州科技大学}
\date{\today}
\titlegraphic{\includegraphics[height=1.25cm]{usts-logo.jpg}}

\AtBeginSection[]{
  \begin{frame}[plain]
    \sectionpage
  \end{frame}
}

\begin{document}

\begin{frame}
  \titlepage
\end{frame}

\begin{frame}{目录}
  \tableofcontents
\end{frame}

\section{基础版式}

\begin{frame}{清晰的信息层级}
  \begin{itemize}
    \item 使用学校蓝作为标题、项目符号与强调色。
    \item 正文保持白底黑字，兼顾投影环境下的可读性。
    \item 页脚集中呈现作者、报告简称、日期与页码。
  \end{itemize}

  \begin{block}{使用建议}
    每页围绕一个核心观点组织内容；文字较多时优先拆页，而不是继续缩小字号。
  \end{block}
\end{frame}

\begin{frame}{三类信息块}
  \begin{columns}[T,onlytextwidth]
    \column{0.30\textwidth}
      \begin{block}{普通信息}
        适合定义、背景和关键说明。
      \end{block}

    \column{0.30\textwidth}
      \begin{exampleblock}{示例或结论}
        用较亮的学校蓝突出正向信息。
      \end{exampleblock}

    \column{0.30\textwidth}
      \begin{alertblock}{重点提醒}
        用深蓝强化需要特别注意的内容。
      \end{alertblock}
  \end{columns}
\end{frame}

\section{图片与布局}

\begin{frame}[fragile]{插入单张图片并添加图注}
  \begin{columns}[T,onlytextwidth]
    \column{0.46\textwidth}
      \begin{figure}
        \centering
        \includegraphics[width=0.78\linewidth,keepaspectratio]{example-image-a}
        \caption{单张图片示例}
      \end{figure}

      \begin{itemize}
        \item \texttt{width} 控制相对宽度；
        \item \texttt{keepaspectratio} 防止图片变形。
      \end{itemize}

    \column{0.50\textwidth}
      \begin{lstlisting}[language={[LaTeX]TeX},basicstyle=\ttfamily\scriptsize]
\begin{figure}
  \centering
  \includegraphics[
    width=.78\linewidth,
    keepaspectratio
  ]{image-file}
  \caption{图片说明}
\end{figure}
      \end{lstlisting}
  \end{columns}
\end{frame}

\begin{frame}[fragile]{图文双栏与图片裁剪}
  \begin{columns}[T,onlytextwidth]
    \column{0.48\textwidth}
      \begin{block}{适合说明型页面}
        左侧放观点、步骤或数据，右侧放与当前结论直接相关的图片。
      \end{block}

      \begin{lstlisting}[language={[LaTeX]TeX},basicstyle=\ttfamily\scriptsize]
\begin{columns}
  \column{.48\textwidth}
  Text or bullets
  \column{.48\textwidth}
  \includegraphics[
    width=\linewidth,
    trim=18 18 18 18,
    clip
  ]{image-file}
\end{columns}
      \end{lstlisting}

    \column{0.46\textwidth}
      \centering
      \includegraphics[
        width=\linewidth,
        trim=18 18 18 18,
        clip
      ]{example-image-b}

      \vspace{0.6em}
      {\small \texttt{trim = left bottom right top}，并与 \texttt{clip} 配合。}
  \end{columns}
\end{frame}

\section{常用内容}

\begin{frame}{公式、表格与双栏内容}
  \begin{columns}[T,onlytextwidth]
    \column{0.47\textwidth}
      \begin{block}{公式示例}
        对于样本均值，常用表达为
        \[
          \bar{x}=\frac{1}{n}\sum_{i=1}^{n}x_i.
        \]
      \end{block}

      \begin{enumerate}
        \item 提出问题；
        \item 展示方法；
        \item 总结结论。
      \end{enumerate}

    \column{0.47\textwidth}
      \begin{table}
        \centering
        \caption{示例对比}
        \begin{tabular}{lcc}
          \toprule
          方案 & 指标 A & 指标 B \\
          \midrule
          基线 & 72 & 0.81 \\
          改进 & 86 & 0.90 \\
          \bottomrule
        \end{tabular}
      \end{table}
  \end{columns}
\end{frame}

\begin{frame}{定义、定理与证明环境}
  \begin{columns}[T,onlytextwidth]
    \column{0.43\textwidth}
      \begin{definition}[算术平均数]
        对于 $n$ 个实数 $x_1,\ldots,x_n$，算术平均数定义为
        \[
          \bar{x}=\frac{1}{n}\sum_{i=1}^{n}x_i.
        \]
      \end{definition}

    \column{0.53\textwidth}
      \begin{theorem}
        \normalfont
        对任意正整数 $n$，有
        \[
          1+2+\cdots+n=\frac{n(n+1)}{2}.
        \]
      \end{theorem}

      \begin{proof}
        \normalfont
        将数列首尾配对，每一对之和均为 $n+1$，整理即可得到结论。
      \end{proof}
  \end{columns}
\end{frame}

\begin{frame}[fragile]{用覆盖层逐步显示内容}
  \begin{columns}[T,onlytextwidth]
    \column{0.43\textwidth}
      \begin{enumerate}[<+->]
        \item 先说明问题背景；
        \item 再展示分析过程；
        \item 最后给出关键结论。
      \end{enumerate}

      \uncover<4->{
        \begin{exampleblock}{呈现原则}
          逐步显示应服务于讲解节奏，而不是制造无意义的动画。
        \end{exampleblock}
      }

    \column{0.53\textwidth}
      \begin{lstlisting}[language={[LaTeX]TeX},basicstyle=\ttfamily\scriptsize]
\begin{enumerate}[<+->]
  \item First point
  \item Second point
  \item Final point
\end{enumerate}

\uncover<4->{Final result}
      \end{lstlisting}
  \end{columns}
\end{frame}

\begin{frame}[fragile]{代码与技术内容}
  \begin{lstlisting}[language=Python]
def mean(values):
    """Return the arithmetic mean."""
    return sum(values) / len(values)
  \end{lstlisting}

  \begin{itemize}
    \item 模板仅依赖常见的 TeX Live 宏包。
    \item 中文文档建议使用 XeLaTeX 编译。
  \end{itemize}
\end{frame}

\begin{frame}[fragile]{超链接、按钮与页内跳转}
  \begin{columns}[T,onlytextwidth]
    \column{0.44\textwidth}
      \begin{block}{外部链接}
        \href{https://www.usts.edu.cn/}{\beamergotobutton{苏州科技大学官网}}
      \end{block}

      \begin{block}{页内跳转}
        \hyperlink{summary}{\beamergotobutton{跳到总结页}}
      \end{block}

    \column{0.52\textwidth}
      \begin{lstlisting}[language={[LaTeX]TeX},basicstyle=\ttfamily\scriptsize]
\href{https://example.com}{
  \beamergotobutton{Website}
}

\hyperlink{summary}{
  \beamergotobutton{Go to summary}
}
      \end{lstlisting}
  \end{columns}
\end{frame}

\begin{frame}[label=summary]{从示例开始制作你的报告}
  \begin{block}{下一步}
    复制 \texttt{example.tex}，替换标题、作者、学院、正文和图片，即可开始制作自己的报告。
  \end{block}

  \begin{itemize}
    \item 保留主题文件，按需要选择页脚模式；
    \item 将图片集中放入 \texttt{assets/}，并使用相对路径引用；
    \item 编译后逐页检查文字换行、图片清晰度和内容密度。
  \end{itemize}
\end{frame}

\end{document}
